% !TeX root = ../main.tex

\chapter{绪论}
\label{cha:introduction}

\section{研究背景与意义}
\label{sec:background}

随着移动互联网与数字内容产业的快速发展，用户面对的信息规模呈指数级增长，信息过载问题日益突出。推荐系统作为连接用户与内容的核心组件，通过建模用户的历史交互行为来预测其潜在兴趣，已广泛应用于电子商务、内容分发、社交媒体和数字广告等场景，成为提升用户体验与平台转化效率的关键技术\cite{zhang2019deep, li2024recent}。

推荐算法的发展经历了若干重要阶段。早期方法主要依赖逻辑回归、矩阵分解和因子分解机等浅层模型，侧重于建模低阶特征交互\cite{koren2009mf, rendle2010fm}。深度学习的引入使模型具备高阶特征交互与端到端表征学习的能力，Wide\&Deep、DeepFM 等模型显著提升了点击率预测的精度\cite{cheng2016wide, guo2017deepfm}。随后，以 DIN、DIEN 为代表的兴趣建模方法通过注意力机制自适应地刻画用户兴趣的动态演化\cite{zhou2018din, zhou2019dien}，而双塔架构则以向量化检索的方式支撑了大规模在线服务\cite{huang2013dssm}。

然而，上述方法普遍受限于封闭的训练数据，只能从用户—物品交互中挖掘协同信号，缺乏对交互背后更广泛世界知识的利用。为突破这一瓶颈，研究者先后引入两类外部知识：一是以知识图谱为代表的结构化知识，通过实体与关系的显式建模增强推荐的可解释性与归纳能力\cite{wang2019kgat, wang2021kgin}；二是以大语言模型为代表的语义知识，凭借其强大的语言理解与开放世界知识直接生成推荐结果或增强既有架构\cite{zhao2024llmrec, wu2024llm4rec}。

尽管两条路线各自取得了显著进展，但将二者结合仍面临若干关键挑战。在知识表示层面，大语言模型擅长处理非结构化文本，知识图谱侧重于结构化关系建模，二者的表示空间存在根本性差异，简单的特征拼接难以充分发挥其互补优势。在知识融合层面，现有方法多将外部知识作为静态特征一次性注入，缺乏对知识质量、相关性与动态演化的考量。在工程落地层面，大语言模型的推理开销与推荐系统对低时延、高吞吐的要求之间存在尖锐矛盾。

针对上述问题，本文从两个互补的视角研究知识增强的推荐方法：其一，将知识以大语言模型生成的语义向量形式注入推荐，解决表示对齐与推理效率问题；其二，将知识以多视角图结构的形式注入推荐，解决图结构动态构建与长尾分布问题。研究成果已在媒体内容推荐场景中完成工程验证，对提升推荐系统的准确性与可部署性具有理论意义与应用价值。

% TODO: 在此补充工程背景（脱敏表述），例如"本文研究工作依托某省级媒体内容平台的智能推荐升级项目"，
%       以及广西重点研发计划（No. 2024AB08144）中跨域推荐中台的应用背景。注意不要出现企业名称与合同金额。

\section{国内外研究现状}
\label{sec:related-work-brief}

围绕知识增强的推荐方法，国内外已形成大语言模型增强推荐与知识图谱增强推荐两条主要技术路线，并在图采样、多智能体协同等方向出现了新的研究趋势。相关工作的详细梳理与对比分析见第~\ref{cha:related-work}~章。

\section{本文主要研究内容}
\label{sec:content}

本文围绕“大语言模型如何有效服务于推荐”这一核心问题，分别从兴趣建模与知识图谱两个落点展开研究。两种方法的共同点在于均以大语言模型作为外部知识源，区别在于大语言模型承担的角色不同：前者作为知识编码器，将开放域知识压缩为可复用的语义画像；后者作为图构建器，将开放域知识组织为显式的图结构。具体研究内容如下。

\subsection{大语言模型增强的用户兴趣聚合表征推荐方法}
\label{subsec:content-llm}

针对现有方法知识注入方式单一、兴趣建模粗粒度且推理开销大的问题，本文提出一种以大语言模型为知识编码器、以聚类兴趣原型为聚合中心的推荐框架。该方法首先利用大语言模型生成物品静态语义画像与用户、物品的动态画像；其次对物品静态画像向量聚类得到兴趣原型中心，并以目标物品的兴趣原型为查询，通过协同过滤特定注意力、物品动态注意力与用户动态注意力聚合用户兴趣表征；最后引入语义难负例对比学习对齐协同过滤嵌入与语义表示，并通过离线预计算与在线内积实现高效推理。相关内容见第~\ref{cha:llm-interest}~章。

\subsection{基于多智能体图构建与图采样的推荐方法}
\label{subsec:content-graph}

针对知识图谱构建依赖人工、图结构静态且长尾节点表征学习不充分的问题，本文提出一种以大语言模型为图构建器、多智能体协同的多视角图推荐框架。该方法由知识融合智能体与语义连接智能体分别生成知识连接与语义连接，与交互连接共同构成多视角图；由长尾增强智能体对稀疏用户补全伪链路、由头部去噪智能体筛除噪声交互；在此基础上构建交互、知识、语义三个视角的子图执行图卷积并融合各视角表征；最终支持直接模式与智能体模式两种推荐方式。相关内容见第~\ref{cha:multi-agent-graph}~章。

\section{本文主要贡献}
\label{sec:contribution}

本文的主要贡献可归纳为以下四个方面：

\begin{enumerate}
  \item 提出了一种以大语言模型为知识编码器的用户兴趣聚合表征推荐框架，将开放域语义知识压缩为可复用的画像向量与兴趣原型，突破了传统协同过滤的知识瓶颈。
  \item 设计了以聚类兴趣原型为中心的多种注意力聚合机制，在保持双塔模型推理效率的同时获得了细粒度兴趣建模的精度。
  \item 提出了一种以大语言模型为图构建器的多智能体图构建与采样优化方法，实现了知识连接与语义连接的自动化生成，摆脱了对人工知识图谱的依赖，并通过长尾伪链路补全与头部强链路筛选提升了稀疏节点的表征质量。
  \item 在媒体内容推荐场景中完成了工程实现与线上验证，验证了所提方法在真实业务约束下的有效性与可部署性。
\end{enumerate}

\section{论文组织结构}
\label{sec:arrangement}

本文共分为五章，各章内容安排如下：

第一章为绪论，介绍研究背景与意义、国内外研究现状、本文的主要研究内容与贡献，以及全文的组织结构。

第二章为国内外研究现状，从推荐算法演进、大语言模型增强推荐、知识图谱增强推荐以及图采样与多智能体方法四个方面梳理相关工作，并分析现有方法的不足。

第三章为大语言模型增强的用户兴趣聚合表征推荐方法，详细介绍画像生成、兴趣原型表示、多路注意力聚合、语义难负例对比学习与在线推理，并通过实验验证方法的有效性。

第四章为基于多智能体图构建与图采样的推荐方法，详细介绍多视角图构建、图采样优化、多视角图卷积与双模式推荐，并给出实验结果与分析。

第五章为总结与展望，总结全文工作，分析研究不足并展望未来研究方向。
